% =====================================================================
%  1bat-ud2-4.tex  —  HORA 4: L'avaluador d'expedients (valor numèric)
%  (sense preàmbul: s'inclou des de main.tex)
% =====================================================================
\horatitol{Sessió 4 — L'avaluador d'expedients}%
{Calcular el valor numèric d'un polinomi substituint la variable, amb la jerarquia d'operacions.}

\subsection{Idea clau}

El \textbf{valor numèric} d'un polinomi és el resultat de substituir la variable
per un nombre concret i operar. És el que fa, per exemple, un tècnic que aplica
una fórmula de copagament a la renda d'una família.

\begin{idea}
\textbf{Jerarquia d'operacions} (l'ordre importa!):
\begin{enumerate}[leftmargin=1.6em, itemsep=2pt]
  \item Primer, les \textbf{potències} (els quadrats).
  \item Després, els \textbf{productes} i les divisions.
  \item Finalment, les \textbf{sumes} i les \textbf{restes}.
\end{enumerate}
\textbf{Compte amb els signes}, sobretot amb els termes negatius.
\end{idea}

\subsection{Exemples}

\begin{exemple}[valor numèric senzill]
Per al polinomi $P = 3x + 5$, el valor numèric en $x=4$ és
\[
P(4) = 3\per 4 + 5 = 12 + 5 = 17.
\]
\end{exemple}

\begin{exemple}[amb un quadrat]
Per a $C(x) = \num{0,05}\,x^2 - 10x + 500$, calculem el valor en $x=100$:
\[
C(100) = \num{0,05}\per 100^2 - 10\per 100 + 500
       = \num{0,05}\per \num{10000} - \num{1000} + 500.
\]
Fem primer la potència ($100^2=\num{10000}$), després els productes
($\num{0,05}\per\num{10000}=500$ i $10\per 100=\num{1000}$) i, al final, sumes i restes:
\[
C(100) = 500 - \num{1000} + 500 = 0.
\]
\end{exemple}

\subsection{Exercicis resolts}

\begin{exresolt}[l'expedient d'una família]
La fórmula de copagament d'un servei és
\[
C(x) = \num{0,05}\,x^2 - 10x + 500,
\]
on $x$ és la renda mensual disponible (en euros) i $C(x)$ el que paga la família.
Calcula el copagament d'una família amb renda $x=200$.
\solucio
Substituïm $x=200$ i respectem la jerarquia:
\begin{align*}
C(200) &= \num{0,05}\per 200^2 - 10\per 200 + 500 \\
       &= \num{0,05}\per \num{40000} - \num{2000} + 500 &&\text{(primer la potència)}\\
       &= \num{2000} - \num{2000} + 500 &&\text{(després els productes)}\\
       &= 500. &&\text{(sumes i restes)}
\end{align*}
\resposta{La família amb renda $200\eur$ paga un copagament de $500\eur$.}
\end{exresolt}

\begin{exresolt}[comparar dos expedients]
Amb la mateixa fórmula $C(x)=\num{0,05}x^2-10x+500$, calcula i compara el
copagament de dues famílies: una amb $x=50$ i una altra amb $x=150$.
\solucio
Família A ($x=50$):
\[
C(50)=\num{0,05}\per 2500 - 10\per 50 + 500 = 125 - 500 + 500 = 125.
\]
Família B ($x=150$):
\[
C(150)=\num{0,05}\per \num{22500} - 10\per 150 + 500 = \num{1125} - \num{1500} + 500 = 125.
\]
Totes dues paguen el mateix copagament, tot i tenir rendes diferents: un resultat
que convida a discutir si la fórmula és justa (treball d'interpretació de la sessió).
\resposta{$C(50)=125\eur$ i $C(150)=125\eur$: paguen el mateix.}
\end{exresolt}

\subsection{Exercicis per fer}

\begin{exercicis}
\begin{exlist}
  \item Calcula el valor numèric de $P = 4x - 7$ per a:
        \quad (a) $x=3$ \quad (b) $x=0$ \quad (c) $x=10$.
  \item Calcula el valor numèric de $Q = x^2 + 2x - 5$ per a:
        \quad (a) $x=4$ \quad (b) $x=1$ \quad (c) $x=0$.
  \item Amb la fórmula de copagament $C(x)=\num{0,05}x^2-10x+500$, calcula el que
        paga cada expedient i ordena'ls de menys a més:
    \begin{center}
    \begin{tabular}{c c}
    \toprule
    Família & Renda $x$ ($\eur$) \\
    \midrule
    A & 0 \\
    B & 100 \\
    C & 250 \\
    D & 300 \\
    \bottomrule
    \end{tabular}
    \end{center}
  \item Un casal modelitza el cost per infant amb $G(x) = \num{0,2}x^2 - 4x + 60$,
        on $x$ són les hores d'activitat. Calcula $G(5)$ i $G(10)$.
  \item \textbf{(Compte amb el signe)} Calcula el valor numèric de
        $R = -2x^2 + 3x + 1$ per a $x=2$. Indica en quin pas cal anar amb cura
        amb el signe.
  \item \textbf{(Interpretació)} Amb $C(x)=\num{0,05}x^2-10x+500$, calcula
        $C(100)$. El resultat és $0$: explica, en el context del copagament,
        què vol dir que una família pagui $0\eur$.
\end{exlist}
\end{exercicis}

% ---------------------------------------------------------------------
\fullsolucions{Sessió 4}
\begin{solucions}
\begin{sollist}
  \item (a) $5$ \quad (b) $-7$ \quad (c) $33$
  \item (a) $19$ \quad (b) $-2$ \quad (c) $-5$
  \item $A=500\eur$, \;$B=0\eur$, \;$C=\num{1125}\eur$, \;$D=\num{2000}\eur$. \;
        Ordre de menys a més: $B<A<C<D$.
  \item $G(5)=45$; \quad $G(10)=40$
  \item $R=-1$ \;(cal anar amb compte amb el signe del terme $-2x^2$:
        $-2\per 2^2=-8$)
  \item $C(100)=0$: la família no paga res. La fórmula assoleix el seu valor
        \emph{mínim} ($0\eur$) precisament a $x=100$, i torna a créixer tant per
        sota com per sobre d'aquesta renda (per exemple, $C(50)=125$ i $C(200)=500$).
        A la sessió 11 veurem que $x=100$ és l'única arrel del polinomi (la gràfica
        només toca l'eix en aquest punt) i que el copagament mai surt negatiu.
\end{sollist}
\end{solucions}
